\newpage
\pagestyle{fancy}
\lhead{\textsc{Conclusion Générale}}
\renewcommand{\headrulewidth}{0.5pt}
\renewcommand{\footrulewidth}{0.5pt}
\chapter*{Conclusion Générale}
\mtcaddchapter

Ce projet de fin d'études a abouti à la conception et au développement d'une application web de
gestion du service après-vente, spécialisée dans la maintenance des climatiseurs et des systèmes
de surpression.

L'application offre une solution complète et intégrée pour digitaliser les processus SAV :

\textbf{Un référentiel métier structuré} : catalogue d'équipements indépendant des clients,
relation \texttt{ClientEquipement} pour matérialiser les installations, gestion des contrats avec
génération automatique du planning préventif.

\textbf{Une gestion opérationnelle complète} : cycle de vie des interventions préventives et
curatives, déclaration de pannes avec pièces jointes, conversion panne $\rightarrow$ intervention,
vérification de disponibilité des techniciens.

\textbf{Un suivi financier automatisé} : facturation conditionnelle (curative + réalisée + hors
contrat), calcul automatique de la TVA à 19\,\%, gestion des statuts de paiement.

\textbf{Des espaces différenciés par rôle} : tableau de bord global pour l'administrateur, tableau
de bord personnel pour le technicien, espace client dédié.

Sur le plan technique, ce projet a permis de maîtriser une architecture moderne en trois couches
(Frontend Next.js/React, Backend API, Base de données Prisma/MySQL) avec Prisma ORM garantissant
la séparation des préoccupations et l'intégrité transactionnelle des données.

La méthodologie Scrum adoptée a permis un développement itératif et incrémental, avec trois sprints
progressifs couvrant l'ensemble du périmètre fonctionnel.

Ce projet constitue une base solide pour une mise en production future, avec des perspectives
d'évolution comme l'intégration de notifications temps réel, la génération de rapports d'activité
avancés ou l'extension à d'autres types d'équipements de maintenance.
